%!TEX root = ../manuscript.tex

\begin{lemma}[Reduction-Preservation under Single Moves]\label{lem:rescue}
Let $w$ be a Gauss word and suppose $w \to^{*}_{\mathcal{B}} \emptyset$
via some bigon-move sequence. Let $\mu$ be any bigon move available
at $w$. Then $\mu(w)$ also admits a bigon reduction to $\emptyset$.
\end{lemma}

\begin{proof}
By strong induction on $\mathrm{cr}(w)$. The base case $\mathrm{cr}(w)
= 0$ is vacuous: there are no available moves. Assume the lemma
holds for all Gauss words of crossing number strictly less than $n$,
and let $w$ have crossing number $n \geq 1$.

Let $\sigma = (\sigma_1, \sigma_2, \dots, \sigma_m)$ be a bigon
reduction of $w$ to $\emptyset$, and let $\mu$ be a bigon move
available at $w$ with chord set $C_\mu \subseteq \{1, \dots, n\}$.

Let $i \in \{1, \dots, m\}$ be the minimal index such that $\sigma_i$
has chord set intersecting $C_\mu$. (Such $i$ exists because $C_\mu$
is eventually removed by $\sigma$.) For every $j < i$, the chord set
of $\sigma_j$ is disjoint from $C_\mu$.

\paragraph{Step 1 (commutation).}
We show that for every $j \in \{0, \dots, i-1\}$, $\mu$ is available
at $w^{(j)} := \sigma_j(\sigma_{j-1}(\cdots \sigma_1(w) \cdots))$
(with $w^{(0)} = w$).

The base case $j = 0$ is the hypothesis. For the induction step,
assume $\mu$ is available at $w^{(j-1)}$. The move $\sigma_j$ has
chord set disjoint from $C_\mu$ (since $j-1 < i$), and $\sigma_j$ is
a bigon move at $w^{(j-1)}$. By Lemma~\ref{lem:del-mon}, $\mu$
remains available at $\sigma_j(w^{(j-1)}) = w^{(j)}$ and $\mu$ and
$\sigma_j$ commute at $w^{(j-1)}$. Hence $\mu$ is available at
$w^{(i-1)}$.

Iterated commutation also yields
\begin{equation}\label{eq:commute}
\mu(w^{(i-1)}) = \sigma_{i-1}(\sigma_{i-2}(\cdots \sigma_1(\mu(w))
\cdots)),
\end{equation}
and each $\sigma_j$ (for $j < i$) is available at the corresponding
state with $\mu$ applied first.

\paragraph{Step 2 (shape of $\sigma_i$).}
Since $\sigma_i$ has chord set intersecting $C_\mu = \{a\}$ (if
$\mu = \text{R1-down}(a)$) or $C_\mu = \{a, b\}$ (if $\mu =
\text{R2-down}(a, b)$), and $\sigma_i$ is a bigon move at
$w^{(i-1)}$, we enumerate the possibilities using
Lemmas~\ref{lem:E}~and~\ref{lem:del-mon}.

\emph{Case A:} $\mu = \text{R1-down}(a)$.

At $w$, $a$'s tokens are cyclically adjacent, so the cyclic interval
$(a_1, a_2)$ contains no endpoints of other chords. The earlier
moves $\sigma_1, \dots, \sigma_{i-1}$ do not touch $a$, and they can
only delete tokens elsewhere, so at $w^{(i-1)}$ the interval
$(a_1, a_2)$ is still empty. Hence $a$ directly nests no chord at
$w^{(i-1)}$, and $\text{R2-down}(a, e)$ is not available. The only
bigon move involving $a$ at $w^{(i-1)}$ is $\text{R1-down}(a) = \mu$.
Hence $\sigma_i = \mu$.

\emph{Case B:} $\mu = \text{R2-down}(a, b)$ with $a$ directly nesting
$b$ at $w$.

At $w$, the bigon condition for $(a, b)$ gives: the outer arcs of
$(a, b)$ contain no endpoints of other chords. Earlier moves do not
touch $a$ or $b$, and can only delete tokens, so at $w^{(i-1)}$ these
outer arcs remain empty. The interior $(b_1, b_2)$ may however have
been affected: at $w$, it can contain chord endpoints, and earlier
moves can remove some.

The bigon moves at $w^{(i-1)}$ involving $a$ or $b$ are exactly:
\begin{itemize}
\item[(B.i)] $\text{R1-down}(a)$, requiring $a$'s tokens adjacent at
  $w^{(i-1)}$;
\item[(B.ii)] $\text{R1-down}(b)$, requiring $b$'s tokens adjacent at
  $w^{(i-1)}$;
\item[(B.iii)] $\text{R2-down}(a, e)$, requiring $a$ directly nests
  $e$ and the bigon condition for $(a, e)$;
\item[(B.iv)] $\text{R2-down}(f, b)$, requiring $f$ directly nests
  $b$ and bigon condition;
\item[(B.v)] $\text{R2-down}(a, b) = \mu$.
\end{itemize}

We analyze each in turn.

(B.i) is impossible: $b$'s tokens lie between $a$'s tokens at $w$
(since $a$ nests $b$), and earlier moves do not touch $b$, so $b$'s
tokens are still between $a$'s at $w^{(i-1)}$. Hence $a$'s tokens are
not adjacent.

(B.iii) is impossible: at $w^{(i-1)}$, the chord directly nested by
$a$ is $b$ (since $a$ nests $b$ at $w^{(i-1)}$, by the above, and no
other chord can lie between them: such a chord would need to have
both tokens in the outer arcs of $(a, b)$ or in the interior
$(b_1, b_2)$; the outer arcs are empty, and if the chord were entirely
inside $(b_1, b_2)$ it would be nested by $b$, not between $a$ and
$b$). By Lemma~\ref{lem:E}, $a$'s direct descendant at $w^{(i-1)}$ is
unique, so $a$ directly nests only $b$, and $\text{R2-down}(a, e)$
with $e \neq b$ is unavailable.

(B.iv) is impossible by a symmetric argument: at $w^{(i-1)}$, $b$ is
directly nested by $a$ and by $a$ alone (Lemma~\ref{lem:E} applied to
$b$'s unique parent in the nesting forest of $w^{(i-1)}$).

That leaves (B.ii) or (B.v). In the (B.v) sub-sub-case,
$\sigma_i = \mu$, and we proceed as in Step~3 below.

In the (B.ii) sub-sub-case, $\sigma_i = \text{R1-down}(b)$, so
$b$'s tokens are adjacent at $w^{(i-1)}$. We handle this via
induction: let $y := \sigma_i(w^{(i-1)}) = w^{(i-1)} \setminus \{b\}$,
which has crossing number $n - k < n$ where $k$ is the number of
distinct chord labels removed by $\sigma_1, \dots, \sigma_i$ (at least
$1$, since $\sigma_i$ removes $b$).

At $y$, the chord $a$'s tokens are now cyclically adjacent: before
$\sigma_i$, they were separated only by $b$'s two tokens (the outer
arcs of $(a, b)$ being empty, and $b$ directly nested by $a$ with no
other chords between them); after removing $b$'s two tokens,
$a$'s tokens become adjacent. Hence $\text{R1-down}(a)$ is available
at $y$.

The remaining sequence $(\sigma_{i+1}, \dots, \sigma_m)$ reduces $y$
to $\emptyset$. By the induction hypothesis applied to $y$ (which has
smaller crossing number than $w$) with move $\text{R1-down}(a)$,
$\text{R1-down}(a)(y)$ also admits a bigon reduction to $\emptyset$.
But $\text{R1-down}(a)(y) = y \setminus \{a\} = w^{(i-1)} \setminus
\{a, b\} = \mu(w^{(i-1)})$.

By Equation~\eqref{eq:commute} and the availability of each
$\sigma_j$ for $j < i$ at $\mu(w), \sigma_1(\mu(w)), \dots$, applying
$\sigma_1, \sigma_2, \dots, \sigma_{i-1}$ to $\mu(w)$ yields
$\mu(w^{(i-1)})$. Prepending these to a bigon reduction of
$\mu(w^{(i-1)})$ to $\emptyset$ yields a bigon reduction of $\mu(w)$
to $\emptyset$.

\paragraph{Step 3 (when $\sigma_i = \mu$).}
If $\sigma_i = \mu$ (in Case A, or in Case B sub-sub-case (B.v)),
apply Equation~\eqref{eq:commute}: starting from $\mu(w)$, the
sequence $\sigma_1, \sigma_2, \dots, \sigma_{i-1}$ (each available by
Step~1) yields $\mu(w^{(i-1)})$. Since $\sigma_i = \mu$ applied at
$w^{(i-1)}$ equals $\mu(w^{(i-1)})$, the original sequence $\sigma$
from $w^{(i-1)}$ onward reads $\sigma_i(w^{(i-1)}) = \mu(w^{(i-1)})
\to^{*} \emptyset$ via $\sigma_{i+1}, \dots, \sigma_m$. Hence
$(\sigma_1, \dots, \sigma_{i-1}, \sigma_{i+1}, \dots, \sigma_m)$
reduces $\mu(w)$ to $\emptyset$.

\paragraph{Conclusion.}
In every case, $\mu(w)$ admits a bigon reduction to $\emptyset$. The
induction step is complete.
\end{proof}
